% Pursuing Stacks (cotes 134-2 to 134-8): archive cote 134-3, site page 6 (PDF page 7), typescript page 6.
% typescript. Draft compared: ps2.tex:241-296.
% AI-transcribed by Claude Opus 5.5 (claude-opus-5-5, adaptive thinking, effort high), Message Batches API, batch ps1 (msgbatch_01Rv5hCaNTtUt4JHx1yqoHPL), 2026-11-12; not yet checked by eye.
\scanpage{134-3/6}
\margin{6}
\mspage{6}
restriction on the coefficient stacks under consideration.

\xout{Workingxoutxaxtheoryxof} Technically speaking, $\infty$-stacks are the
common denominator of n-stacks for arbitrary n, in much the same way as
n-stacks appear both as the next-step generalization of(n-1)-stacks
(\inserted[typed]{the former}forming a category which admit\hand{s} the category of (n-1) stacks as a full
subcategory), and as the most natural "higher order structure" appearin
% NOTE: the right edge of the scan cuts off the end of this line ("appearin...")
on the category of all \xout{nxstacks} (n-1)-stacks, \inserted{(and on} \xout{xxxx} various \unsure{analoguous|analogous} ca-
tegories whose objects are (n-1)-stacks subject to some restrictions or
endowed with some extra structure\hand{)}. I'll have to make this explicit in
due course. For the time being, when speaking of ``stacks'' or ``Gr-stacks\hand{''}
it will be understood (unless otherwise specified) that we are dealing
with the infinite order notions, which encompass\hand{\unsure{es}} the finite ones.

Working out a theory of stacks over topoi, as the natural founda-%
tions of non commutative cohomological algebra, would amount among
others to rewrite Giraud's book within this considerabl\over{e}{y} wider framework
% NOTE: the right edge of the scan cuts off the end of this line after "framework"
Of course, this mere prospect would'nt be particularly exciting by it-%
self, if it did not appear as something more than grinding through an
unending exercise of rephrasing and reproving known things, replacing
everywhere n=0,1 or 2 by arbitrary n. I am convinced however that there
is a lot more to it - namely the fascination of gradually discovering
and naming and getting acquainted with presently still unknown, unnamed,
mysterious structures. As is the case so often when making a big step
backwards for gaining new perspective, there is not merely a quantitati
ve change (from n\hand{$\leqslant$}2 to arbitrary n say), but a qualitative change in
scope and depth of vision. One such step was already taken I feel by
Daniel Quillen and others, when realizing that homotopy constructions
make sense not only in the usual homotopy category, or \inserted{in} one or the other
categories of models which give rise to it, but in more or less arbitra
\margin{-}ry categories, by working with semi-simpli\over{\unsure{ci}}{ci}al objects \over{\unsure{of}}{in} these. The
step I am proposing is of a \inserted{somewhat} different type. \uline{The notion of a stack}
\uline{here appears as the unifying} \inserted[typed]{concept}\xout{notion} \uline{for a synthesis of homotopical alge-}%
\uline{bra}, \uline{and non commutative cohomological algebra}. This$_{\text{x}}$ \xout{Ixfxxxx} (rather
than \inserted[typed]{merely} furnishing us with still another description of homotopy types,
more convenient for expression of the homotopy groups) se\over{\unsure{ms}}{ems} to me the
real ``raison d'être'' of the notion of a stack, and the main motivation
for pushing ahead a theory of stacks.
